\StartChapterOnNewColumn{Komplexe Zahlen\ZSFScriptRef{14,18,19}}

\SubsectionBar{Eigenschaften}

Eine \ZSFkeyword{komplexe Zahl} ist von der Form $z = a + ib$ wobei $a = \operatorname{Re}(z)$ der \ZSFkeyword{Realteil} und $b = \operatorname{Im}(z)$ der \ZSFkeyword{Imaginärteil} ist.

Die \ZSFkeyword{komplex konjugierte Zahl} $\bar{z} = a - ib$ ist die Spiegelung von $z$ an der reellen Achse (meist die $x$-Achse).

\textVorBox{Es gelten die folgenden Regeln für Konjugation und Betrag:}

\begin{fulltablebox}{C | C}
    \TableTitle{Konjugation} & \TableTitle{Betrag} \\ \hline
    \ZSFrowColor
    $\displaystyle \overline{z_1 + z_2} = \overline{z_1} + \overline{z_2}$ & $\displaystyle |z_1 z_2| = |z_1||z_2|$ \\ \hline
    $\displaystyle \overline{z_1 z_2} = \overline{z_1} \, \overline{z_2}$ & $\displaystyle \left| \frac{z_1}{z_2} \right| = \frac{|z_1|}{|z_2|}$ \\ \hline
    \ZSFrowColor
    $\displaystyle \overline{z^n} = \bar{z}^n$ & $\displaystyle |z_1 + z_2| \leq |z_1| + |z_2|$ \\ \hline
    \multicolumn{2}{c}{$\displaystyle z\bar{z} = |z|^2$} \\
\end{fulltablebox}

Die Darstellung $z = r(\cos(\varphi) + i\sin(\varphi))$ nennt man \ZSFkeyword{Polarform}.

\textVorBox{Die Darstellung $z = re^{i\varphi}$ nennt man \ZSFkeyword{Eulersche Form}. Es gilt:}
\begin{formulabox}
    $\displaystyle r = |z| = \sqrt{a^2 + b^2}$
\end{formulabox}

\textVorBox{und:}
\begin{fulltablebox}{C | C | C}
    \TableTitle{$a > 0$} & \TableTitle{$a < 0,\ b \geq 0$} & \TableTitle{$a < 0,\ b < 0$} \\ \hline
    \ZSFrowColor
    $\displaystyle \varphi = \arctan\left(\frac{b}{a}\right)$ & $\displaystyle \varphi = \arctan\left(\frac{b}{a}\right) + \pi$ & $\displaystyle \varphi = \arctan\left(\frac{b}{a}\right) - \pi$ \\
\end{fulltablebox}
\textNachBox{wobei $\varphi$ gegen den Uhrzeigersinn gemessen wird.}

\textVorBox{\ZSFdanger{Wichtig:}}
\begin{formulabox}
    $\displaystyle z = e^{i\pi} = -1,\quad z = e^{i\frac{\pi}{2}} = i$
\end{formulabox}

\SubsectionBar{Rechenregeln}

\textVorBox{Die Eulersche Form eignet sich für komplizierte Operationen mit komplexen Zahlen:}

\begin{fulltablebox}{Q{0.7} | f{1.3}}
    \ZSFrowColor Multiplikation: & \displaystyle z_1 \cdot z_2 = r_1 \cdot r_2 \cdot e^{i(\varphi_1 + \varphi_2)} \\ \hline
    Division: & \displaystyle \frac{z_1}{z_2} = \frac{r_1}{r_2} e^{i(\varphi_1 - \varphi_2)} \\ \hline
    \ZSFrowColor Inverse: & \displaystyle \frac{1}{z} = \frac{1}{r} e^{-i\varphi} \\ \hline
    Potenz: & \displaystyle z^n = r^n e^{in\varphi} \\ \hline
    \ZSFrowColor Wurzel: & \displaystyle \sqrt[n]{z} = \sqrt[n]{r} e^{i\left(\frac{\varphi + 2k\pi}{n}\right)} \\
\end{fulltablebox}
\textNachBox{wobei $k = 0, 1, 2, \dots, n-1$ \newline Alle Wurzeln liegen auf Kreis mit Radius $\sqrt[n]{r}$}

\SubsectionBar{Polynome\ZSFScriptRef{21,57}}

\textVorBox{Ein \ZSFkeyword{Polynom} ist eine Summe von Potenzen einer Variable:}
\begin{formulabox}
    $\displaystyle p(x) = ax^n + bx^{n-1} + \dots + cx + d$
\end{formulabox}
\textNachBox{$x$ wird als \ZSFkeyword{Nullstelle} des Polynoms $p$ bezeichnet, wenn $p(x) = 0$.}

Der \ZSFkeyword{Grad} eines Polynoms ist dessen grösster Exponent (oben $n$).

Ein Polynom vom Grad $n$ hat $n$ Nullstellen.

Ein Polynom kann reelle und komplexe Nullstellen besitzen. Diese können einfach ($p(x) = 0$), doppelt ($p(x) = 0$, $p'(x) = 0$), dreifach ($p(x) = 0$, $p'(x) = 0$, $p''(x) = 0$), bis $n$-fach vorkommen.

Komplexe Nullstellen kommen immer paarweise vor ($z, \bar{z}$). Ein Polynom mit ungeradem Grad besitzt mindestens eine reelle Nullstelle. Eine einfach komplexe Nullstelle entspricht zwei Nullstellen. Eine doppelt komplexe Nullstelle entspricht vier Nullstellen.

\textVorBox{Es gelten folgende Regeln für den Grad eines Polynoms:}
\begin{formulabox}
    \begin{longformula}
        \operatorname{Grad}\{p + q\} &\leq \max(\operatorname{Grad}\{p\}, \operatorname{Grad}\{q\}) \\
        \operatorname{Grad}\{pq\} &= \operatorname{Grad}\{p\} + \operatorname{Grad}\{q\}
    \end{longformula}
\end{formulabox}
